\scanpage{14}
\mspage{11}

Si $\mathbb D$ est un dérivateur, la catégorie

\[
\unsure{\mathcal A|\mathscr A}=A_{\mathbb D}=\mathbb D(e)
\]

joue un rôle important. On l'appelle la \uline{catégorie de base}, ou la \uline{cat. fondamentale}, du dérivateur, ses objets sont les coefficients fondamentaux. Pour $I$ quelconque dans $\mathrm{Dia}$, on a donc pour $F\in\mathbb D(I)$ coeff. sur $I$

\[
I^{\unsure{\circ|o}}\longrightarrow\deleted{\unsure{\mathbb D(e)}}\mathcal A\ ,\qquad i\longmapsto F_{i}
\]

un diagramme de type $I^{\circ}$ dans $\mathcal A$, fonctoriellement, i.e.

\[
\mathbb D(I)\longrightarrow\underline{\mathrm{Hom}}\,(I^{\circ},\deleted{\unsure{\mathbb D}}\underset{\substack{\Vert\\\mathbb D(e)}}{\mathcal A})
\]

(cf page 8)

\uline{Der}$_3$ : Pour \unsure{tt} $I\in\mathrm{Dia}$, le foncteur est \uline{conservatif}, i.e. pour que

% slanted
\margin{\uline{NB} Mais ce foncteur n'a aucune raison d'être \uline{fidèle} !}

\[
F\longrightarrow G
\]

dans $\mathbb D(I)$ soit iso, il \unsure{f et s|fds} que $\forall i\in I$,

\[
F_{i}\longrightarrow G_{i}
\]

le soit.

\uline{Cor.} Le foncteur canonique (pour $I,J\in\mathrm{Dia}$)

\[
\mathbb D(I\times J)\overset{\Theta}{\longrightarrow}\underline{\mathrm{Hom}}\,(J^{\circ},\mathbb D(I))
\]

est conservatif. En effet, si $F$ est de $\mathbb D(I\times J)$, on a pour $\unsure{i,j|ij}\in I\times J$

\[
F_{i,j}\simeq\Theta(F)(j)_{i}\ .
\]

% archive sheet number: 13
